\documentclass[runningheads]{llncs}

\usepackage[T1]{fontenc}
\usepackage[utf8]{inputenc}
\usepackage{graphicx}
\usepackage{hyperref}
\usepackage{booktabs}
\usepackage{amsmath}
\usepackage{parskip}
\usepackage{float}

\begin{document}
\raggedbottom
\title{ArtBench-10 Generative Modeling: A Comparative Study of VAE, GAN, and Diffusion Models}
\author{Tiago Silva \and João Coelho}
\institute{University of Coimbra}
\maketitle

\begin{abstract}
    Nothing yet.
\end{abstract}

\section{Introduction}

\section{Problem Formulation}


\section{Methodological Framework}

\subsection{Data Processing Pipeline}

\subsubsection{VAE}

\subsubsection{DCGAN}

\subsubsection{Diffusion Models}

\subsection{Implementation Details and Reproducibility}

\subsection{Theoretical Rationale for Hyperparameter Choices}

\subsubsection{Variational Autoencoders (VAE)}


\subsubsection{Generative Adversarial Networks (GAN)}

\subsubsection{Diffusion Models}


\section{Experimental Design}
\subsection{Hyperparameter Optimization Strategy}
We performed hyperparameter optimization with Bayesian optimization (Optuna). To keep the search feasible on standard hardware, 
tuning was conducted on a development subset comprising 20\% of the training data. For each model family, candidate configurations 
were trained on a training split and selected using performance on a held-out validation split. We minimized the following 
surrogate objectives: (i) \textit{VAE}: validation reconstruction loss; (ii) \textit{GAN}-based models: validation sum of discriminator 
and generator losses; (iii) \textit{Diffusion models}: validation denoising loss under a fixed Gaussian noise schedule.
The hyperparameter search spaces are summarized in Table~\ref{tab:hpo_spaces}, and the best configurations found on the studies are 
reported in Table~\ref{tab:hpo_best}. We used validation losses as surrogates during tuning because FID/KID are computationally expensive 
and high-variance at the per-trial scale: they require generating large sample sets and repeated feature extraction. In contrast, validation losses are 
available at each epoch and exhibit lower variance which improves sample efficiency under limited computational resources. We note, however, that GAN losses are far from beeing perfect proxies 
for perceptual sample quality; therefore, we use this objective only for hyperparameter selection during tuning and rely on qualitative inspection 
for the selected final parameters.
\begin{table}[H]
    \centering
    \caption{Hyperparameter search spaces used in Bayesian optimization.}
    \label{tab:hpo_spaces}
    \begin{tabular}{p{0.2\textwidth}p{0.3\textwidth}p{0.5\textwidth}}
        \toprule
        \textbf{Model} & \textbf{Parameters} & \textbf{Values / Range} \\
        \midrule
        VAE &
        latent\_dim, base\_channels, batch\_size &
        $\{16,32,64\}$; $\{16,32,64\}$; $\{32,64,128\}$ \\
        & lr, $\beta$ &
        $[10^{-4}, 3\times 10^{-3}]$ log-uniform; $[10^{-6}, 10^{-3}]$ log-uniform \\
        \midrule
        DCGAN &
        latent\_dim, feature\_maps, batch\_size &
        $\{128,256,512\}$; $\{32,64\}$; $\{32,64,128\}$ \\
        & lr, $\beta_1$, real\_label &
        $[10^{-4}, 5\times 10^{-4}]$ log-uniform; $[0.3,0.7]$; $[0.85,1.0]$ \\
        \midrule
        PixelUNet &
        model\_channels, batch\_size, num\_timesteps &
        $\{32,64,96\}$; $\{16,32,64\}$; $\{500,1000\}$ \\
        & lr &
        $[10^{-4}, 5\times 10^{-4}]$ log-uniform \\
        \midrule
        LatentDenoiser &
        model\_channels, num\_res\_blocks, batch\_size &
        $\{32,64,96\}$; $\{2,3,4\}$; $\{16,32,64\}$ \\
        & lr, num\_timesteps &
        $[10^{-4}, 5\times 10^{-4}]$ log-uniform; $\{500,1000\}$ \\
        \bottomrule
    \end{tabular}
\end{table}



\begin{table}[H]
    \centering
    \caption{Best configurations obtained.}
    \label{tab:hpo_best}
    \begin{tabular}{p{0.28\textwidth}cccc}
        \toprule
        \textbf{Parameter} & \textbf{VAE} & \textbf{DCGAN} & \textbf{PixelUnet} & \textbf{LatentDenoiser}\\
        \midrule
        Best objective value & 0.02658 & 1.68158 & 0.0358 & 0.2598 \\
        Learning rate & 0.0005317 & 0.0001261 & 0.0004748 & 0.0001769\\
        Batch size & 32 & 32 & 16 & 32 \\
        Latent dimension & 32 & 256 & 64 & 64 \\
        Base channels / feature maps & 64 & 64 & 64 & 64 \\
        $\beta$ (VAE) & $1.14 \times 10^{-6}$ & -- & -- & -- \\
        $\beta_1$ (Adam, DCGAN) & -- & 0.4813 & -- & -- \\
        real\_label (DCGAN) & -- & 0.9836 & -- & -- \\
        num\_timesteps (diffusion) & -- & -- & 1000 & 1000 \\
        num\_res\_blocks (LatentDenoiser) & -- & -- & -- & 4 \\
        \bottomrule
    \end{tabular}
\end{table}


\subsection{Evaluation Protocol}
Across the selected configurations (Table~\ref{tab:hpo_best}), 
the optimization process consistently favored smaller batch sizes. 
This aligns with the well-documented generalization gap observed in 
large-batch training~\cite{keskar}. Specifically, the inherent noise 
in small-batch stochastic gradients serves as a form of implicit 
regularization, facilitating the escape from sharp local minima and 
leading the model toward flatter, more robust regions of the loss 
landscape. Regarding the learning rates, the observed values
are relatively low, consistent with the need for stable convergence and
stability during training, especially for GANs which are known to be sensitive to hyperparameter settings.
Latent dimensions reveal important differences across model families: the VAE favored a much smaller 
latent space (32) compared to the DCGAN (256), which may reflect the different in the 
model's architectures: GANs don't pay an explicity penalty (KL divergence)
for using larger latent spaces, while VAEs are encouraged to find compact representations.
A larger latent space may encourage more diverse samples as long as the generator
actually uses it. For the VAE, the optimal $\beta$ value is very low, 
indicating that the model prioritizes reconstruction fidelity over latent space regularization.
For the DCGAN, the optimal $\beta_1$ value is around 0.48, which is close 
to the commonly used value of 0.5~\cite{dcgan}, suggesting a balance between stability and convergence speed. 
The real\_label parameter being close to 1 (0.9836) indicates that the discriminator 
benefits from a small anti-overconfidence margin, which can help prevent it from becoming too confident 
and overpowering the generator during training without weakening the learning signal too much.

\section{Results and Interpretation}

\subsection{Ablation Studies and Discussion}

\section{Conclusion}

\section*{Use of External Resources and AI Tools}
This project utilized ArtBench-10 (Kaggle). Gemini assisted in LaTeX formatting and analysis verification.

\nocite{artbench,vae,dcgan,fid,keskar}
\bibliographystyle{splncs04}
\bibliography{main}

\end{document}
